\documentclass{beamer}
\usepackage{hyperref}
\usepackage[czech]{babel}

%\usetheme{Boadilla}
\usetheme{default}
\usecolortheme{seahorse}
\definecolor{myblue}{rgb}{0.2, 0.2, 0.6}
\setbeamertemplate{caption}{{\color{myblue}Obrázek:} \raggedright\insertcaption\par}
\setbeamertemplate{footline}[frame number] 

\newcommand{\authorname}{Dluhoš Matěj, Gaja Jan, Gajdošík Petr, Jurásek Petr, Phan Thanh Tú, Jan Rádl}
\newcommand{\authorsshort}{Dluhoš, Gaja, Gajdošík, Jurásek, Phan, Rádl}
\newcommand{\thesisname}{PDS: Python}

\title{\thesisname}
\author{\authorname}
\institute{Univerzita Palackého v Olomouci}
\date{\today}

\setbeamertemplate{navigation symbols}{}
\setbeamertemplate{headline}{}
\setbeamertemplate{footline}{
  \leavevmode%
  \hbox{%
  \begin{beamercolorbox}[wd=.4\paperwidth,ht=2.25ex,dp=1ex,center]{author in head/foot}%
    \usebeamerfont{author in head/foot}\authorsshort
  \end{beamercolorbox}%
  \begin{beamercolorbox}[wd=.3\paperwidth,ht=2.25ex,dp=1ex,center]{title in head/foot}%
    \usebeamerfont{title in head/foot}\thesisname
  \end{beamercolorbox}%
  \begin{beamercolorbox}[wd=.3\paperwidth,ht=2.25ex,dp=1ex,right]{date in head/foot}%
    \usebeamerfont{date in head/foot}\insertshortdate{}\hspace*{2em}
    \insertframenumber{} / \inserttotalframenumber\hspace*{2ex} 
  \end{beamercolorbox}}%
  \vskip0pt%
}

%\logo{\includegraphics[height=1cm]{UP_logo.png}}
\usepackage{tikz}
\addtobeamertemplate{headline}{}{%
    \begin{tikzpicture}[overlay, remember picture]
    	\ifnum\insertframenumber>1
        	\node[anchor=north east, inner sep=5pt] at (current page.north east) {\includegraphics[height=1cm]{obrazky/UP_logo.png}};
        \fi
    \end{tikzpicture}
}

\setbeamercolor{titlelike}{bg=,fg=}

\begin{document}

\begin{frame}
	\titlepage
\end{frame}

\begin{frame}{Obsah}
	\tableofcontents
\end{frame}

\section{Úvod do Pythonu s paralelismem}
\subsection{Historie a zaměření}
\begin{frame}{Historie a zaměření}
	\begin{itemize}
		\item Vytvořen \textbf{Guidem van Rossumem} v roce \textbf{1991}.
	\end{itemize}
	\begin{itemize}
		\item Jednoduchost a čitelnost kódu.
		\item Vysoká produktivita programátorů.
		\item Univerzální použití: od webových aplikací přes vědecké výpočty až po strojové učení.
	\end{itemize}
\end{frame}

\subsection{Paralelismus v Pythonu}
\begin{frame}{Paralelismus v Pythonu}
	\begin{itemize}
		\item Python získal základní podporu paralelismu prostřednictvím vláken ve verzi 1.5.2 (1999).
		\item Zaveden \textbf{GIL} (Global Interpreter Lock)
		\begin{itemize}
			\item [\textendash] Usnadnilo integraci nativního kódu v C.
			\item [\textendash] Vyřešilo problémy s přístupem ke sdíleným datům ve více vláknech.
			\item [\textendash] Optimalizovalo výkon na systémech s jedním jádrem.
			\item [\textendash] Omezuje paralelní výkon na vícejádrových procesorech.
			\item [\textendash] Vede k neefektivnímu využití vláken u \textbf{CPU-bound} úloh.
			\item [\textendash] \hypersetup{urlcolor=blue} \href{https://peps.python.org/pep-0703/}{PEP 703} - Making the Global Interpreter Lock Optional in CPython 
		\end{itemize}
	\end{itemize}
\end{frame}

\subsection{Proč odstranit GIL?}
\begin{frame}{Proč odstranit GIL?}
	\begin{itemize}
		\item Většina zařízení dnes obsahuje vícejádrové procesory, kde GIL představuje významné omezení.
		\item Efektivnější paralelní výpočty, například v oblastech umělé inteligence nebo numerických simulacích.
		\item Výzvy při odstranění GIL:
		\begin{itemize}
			\item [\textendash] Potřeba bezpečného a efektivního spravování sdílené paměti (například více zamykání nebo jiných synchronizačních mechanismů).
			\item [\textendash] Zvýšení složitosti implementace interpretu Pythonu.
		\end{itemize}
	\end{itemize}
\end{frame}

\subsection{Komunikační model Pythonu}
\begin{frame}{Komunikační model Pythonu}

\end{frame}

\section{Synchronizační nástroje standardní knihovny}
\subsection{Knihovna threading}

\begin{frame}{Knihovna threading}
    \begin{itemize}
        \item Vestavený modul pro paralelní spouštění vláken
        \item Omezena GIL
        \item Vlákna běží na stejné haldě, využívají tedy stejný paměťový prostor
        \vskip 0.35in
        \item Často využíváné pro:
        \begin{itemize}
            \item I/O operace (Read/Write, Síťování, user input)
            \item Grafické uživatelské rozhraní 
        \end{itemize}
    \end{itemize}
\end{frame}

\begin{frame}{Třída Thread}
    \begin{itemize}
        \item Instance reprezentuje vlákno
        \item Předání funkce a argumentů
        \vskip 0.2in

        \item \textit{start()} - spustí vytvořené vlákno
        \item \textit{join()} – Počká na dokončení práce vlákna
        \item \textit{is\_alive()} – Pracuje stále vlákno?
        \item \textit{local()} – data konkrétního vlákna

        \vskip 0.2in
        \item Ukázka
        
        \vskip 0.2in
        \item Při práci s vlákny může dojít k synchronizačním problémům
        \item Modul poskytuje mnoho synchronizačních nástrojů
    \end{itemize}
\end{frame}

\begin{frame}{Lock}
    \begin{itemize}
            \item Zajistí, že ke sdíleným zdrojům bude mít přítup pouze jedno vlákno
            \item Můžeme zamknou příkaz nebo skupinu příkazu

            \vskip 0.15in
            \item \textit{acquire()} – zamkne objekt
            \item \textit{release()} – odemkne objekt
            \item klíčové slovo with - zamkne/odemkne blok automaticky

            \vskip 0.15in
            \item Varianta \textbf{RLock} pro rekurzivní funkce
            \item Ukázka
    \end{itemize}
\end{frame}

\begin{frame}{Semaphore}
    \begin{itemize}
        \item Dokáže omezit počet vláken, které mají přístup ke zdroji
        \item Konstruktor očekává value, určuje maximální počet přístupů ke zdroji
        \newline (Varianta Binary Semaphore – maximálně 1 přístup)

        \vskip 0.25in
        \item \textit{acquire()} – zjistí jestli je count $>$ 0, pokud ano, dekrementuje count, jinak zablokuje vlákno
        \item \textit{release()} – inkrementuje count, umožní tak dalšímu vláknu přístup
        \item Opět with statement dělá automaticky
        
        \vskip 0.25in
        \item Ukázka
    \end{itemize}
\end{frame}

\begin{frame}{Barrier}
    \begin{itemize}
        \item Umožňuje zablokovat vlákna, dokud všechna nedojdou na určité místo v programu
        \item Až všechna vlákna dosáhnou bariéry, pokračují dál v práci
        \item Vhodné pokud vlákna pracují ve fázích, kdy po dokončení jedné se spustí další

        \vskip 0.25in
        \item \textit{wait()} – Místo, kde se vlákno zastaví, dokud ostatní vlákna k němu nedojdou
        
        \vskip 0.25in
        \item Ukázka
    \end{itemize}
\end{frame}

\begin{frame}{Event}
    \begin{itemize}
        \item Jednoduchý mechanizmus, který umožňuje signalizaci mezi vlákny
        \item Používá interní příznak, který nabývá hodnoty True nebo False
        \vskip 0.25in

        \item \textit{wait()} – vlákna budou blokována, dokud nebude flag nastaven na True
        \item \textit{set()} – Nastaví flag na True, vlákna volající wait() budou propuštěna
        \item \textit{clear()} – Nastaví flag na False, vlákna volající wait() budou zablokována
        
        \vskip 0.25in
        \item Ukázka
    \end{itemize}
\end{frame}

\begin{frame}{Condition}
    \begin{itemize}
        \item Funguje obdobně, ale používá se současně se zámkem

        \vskip 0.25in
        \item \textit{wait()} – stejné, ale navíc odemkne zámek
        \item \textit{notify()} – odblokuje jedno čekající vlákno
        \item \textit{notify\_all()} – odblokuje všechna čekající vlákna
        
        \vskip 0.25in
        \item Ukázka
    \end{itemize}
\end{frame}

\subsection{Knihovna multiprocessing}
\begin{frame}{Knihovna multiprocessing}
    \begin{itemize}
        \item Modul pro paralelizaci procesů
        \item Nepracuje s vlákny ale s procesy
        \item Umožňuje překonat limitace GIL
        \item Vhodné pro složitější výpočty
        
        \vskip 0.3in
        \item Výhody:
        \begin{itemize}
            \item Každý vytvořený proces běží nezávisle a může být vykonáván na jiných jádrech procesoru
            \item Všechny procesy mají vlastní paměť
            \item Vestavené metody pro bezpečnou práci se sdílenými prostředky
        \end{itemize}
    \end{itemize}
\end{frame}

\begin{frame}{Třída Proces}
    \begin{itemize}
        \item Instance reprezentuje jeden nezávislý proces
        \item Tvorba stejná jako u threading
        \item Stejné funkce
        \vskip 0.25in

        \item Využití Pool
        \begin{itemize}
            \item Umožňuje efektivně spravovat procesy
            \item Automatizuje jejich tvorbu a rozdělení mezi jádra CPU
            \item \textit{map()} - pro paralelní provedení funkce na každý prvek z kolekce dat
        \end{itemize}
        \item Ukázka
    \end{itemize}
\end{frame}

\begin{frame}{Data a komunikace}
    \begin{itemize}
        \item Sdílení dat
        \begin{itemize}
            \item Modul poskytuje objekty Array a Value pro sdílení dat mezi procesy
        \end{itemize}

        \item Komunikace
        \begin{itemize}
            \item Pro plné využití síly procesů je potřeba komunikace mezi nimi
            \item Objekty Queue a Pipe pro zasílání zpráv
            \item S Queue se pracuje stejně jako s klasickou frontou
            \item Pipe je vhodnější pro komunikaci mezi pouze dvěma procesy
        \end{itemize}
        \item Synchronizace
        \begin{itemize}
            \item Stejné nástroje pro synchronizaci jako v modulu threading
            \item Lock, Semaphore, Barrier, Event, Condition
        \end{itemize}
        \item Ukázka
    \end{itemize}
\end{frame}

\subsection{Knihovna asyncio}
\begin{frame}{Knihovna asyncio}
    \begin{itemize}
        \item Umožňuje asynchronní zpracování vstupně-výstupních operací
        \item Založeno na kooperativním multitaskingu
        \item Můžeme se tak vyhnout nečinnosti při čekání na dokončení I/O operací
        \item Vhodné pro síťové operace, práce s databází, čtení souborů…
        
        \vskip 0.35in
        \item \textbf{Event Loop} – koordinuje spuštění asynchronních operací
        \item \textbf{Coroutine} – speciální asynchronní funkce
        \item \textbf{Task} – zajišťuje běh coroutine v event loopu
        \item \textbf{Future} – Předběžný výsledek operace, která nebyla ještě dokončena
    \end{itemize}
\end{frame}

\begin{frame}{Třída Task}
    \begin{itemize}
        \item Při vytváření Tasku předáváme korutinu
        \item Následně je její běh naplánován Event Loopem
        \item Metody pro práci s Tasky během toho co se vykonávají
        \begin{itemize}
            \item \textit{.result()}
            \item \textit{.exception()}
            \item \textit{.cancel()}
        \end{itemize}
        
        \vskip 0.25in
        \item Další užitečné nástroje
        \begin{itemize}
            \item Stejné synchronizační nástroje jako v předchozích modulech

            \item \textit{.wait\_for()} umožňuje nastavit maximální čas čekání na dokončení. Vyvolává TimeoutError
        \end{itemize}
        \item Ukázka
    \end{itemize}
\end{frame}

\section{Nástroje externích knihoven}
\begin{frame}{Nástroje externích knihoven}

\end{frame}

\section{Nástroje pro distribuované systémy}
\subsection{Knihovna socket}
\begin{frame}{Knihovna socket}

\end{frame}

\subsection{Knihovna http.server}
\begin{frame}{Knihovna http.server}

\end{frame}

\subsection{RabbitMQ}
\begin{frame}{RabbitMQ}

\end{frame}

\subsection{Další nástroje}
\begin{frame}{Další nástroje}

\end{frame}

\end{document}